\section{Bounded concentration and finite-class learning: complete proofs}
\label{sec:annex-p15-bounded-concentration-learning}

\begin{annexitem}{annex:p15.conc-exponential}{thm:p15-conc-exponential}{Convex chords, the one-eighth coefficient and exponential Markov}
\begin{proof}
For \(x\in[a,b]\), convexity of the exponential gives
\[
 e^{\lambda x}\leq
       \frac{b-x}{b-a}e^{\lambda a}
            +\frac{x-a}{b-a}e^{\lambda b}.
\]
Taking expectations proves the chord inequality. Set
\(p=(\mu-a)/(b-a)\) and \(s=\lambda(b-a)\).
After centering, the bound is
\[
 Ee^{\lambda(X-\mu)}
       \leq e^{-ps}(1-p+pe^s).
\]
For \(p=0\) or one, nonnegativity of \(X-a\) or \(b-X\)
implies that \(X\) equals the appropriate endpoint almost surely;
the centered moment is one. For \(0<p<1\), let
\(K_p(s)=-ps+\ln(1-p+pe^s)\). Then
\[
 K_p(0)=K'_p(0)=0,\qquad
 K''_p(s)=u_s(1-u_s)\leq1/4,\qquad
 u_s=\frac{pe^s}{1-p+pe^s}\in(0,1).
\]
Taylor's formula with integral remainder, valid also for
negative \(s\), yields \(K_p(s)\leq s^2/8\).
This proves the centered bound. If \(a=b\), the centered
variable is zero and the same formula is an equality.

For a fair binary variable of range one,
\(\ln Ee^{\lambda(X-1/2)}=\ln\cosh(\lambda/2)\).
Its value and first derivative at zero vanish, and its
second derivative at zero is \(1/4\). The ratio of this
logarithm to \(\lambda^2\) tends to \(1/8\).
Any uniform smaller quadratic coefficient would therefore
fail for sufficiently small nonzero \(\lambda\).

Finally, for \(\lambda>0\),
\(\mathbf1_{\{W\geq t\}}\leq e^{\lambda(W-t)}\)
pointwise. Taking its expectation gives the exponential
Markov bound. This argument does not assert equality of
the indicator and the exponential away from their contact.
\end{proof}
\end{annexitem}

\begin{annexitem}{annex:p15.conc-common-range}{thm:p15-conc-common-range}{The optimized bound and the complete mean-dependent refinement}
\begin{proof}
Let \(\mu_i=EX_i\), so \(np=\sum_i\mu_i\).
The chord bound on \([0,1]\), independence and the
arithmetic-geometric mean inequality give, for \(\lambda>0\),
\[
 \begin{split}
 P(S_n\geq nt)
 &\leq e^{-\lambda n(p+t)}
                      \prod_i(1-\mu_i+\mu_i e^\lambda)\\
 &\leq \{e^{-\lambda(p+t)}(1-p+pe^\lambda)\}^n.
 \end{split}
\]
The logarithmic derivative of the last bracket is
\(-(p+t)+pe^\lambda/(1-p+pe^\lambda)\).
Its derivative is strictly positive. Its unique zero is
\[
 \lambda_*=
       \ln\frac{(p+t)(1-p)}{p(1-p-t)}>0.
\]
Substitution gives the bracket
\((p/(p+t))^{p+t}((1-p)/(1-p-t))^{1-p-t}\).
This proves the divergence bound.

To prove the exact refinement, define
\(G(t,p)=D(p+t\Vert p)/t^2\) for \(0<t<1-p\).
Differentiation gives
\[
 t^2\frac{\partial G}{\partial t}
      =H\!\left(\frac{t}{1-p}\right)
           -H\!\left(\frac{t}{p+t}\right),
 \qquad H(z)=(1-2/z)\ln(1-z).
\]
Both arguments lie in \((0,1)\). The convergent logarithmic
series gives
\[
 H(z)=2+\sum_{k=2}^{\infty}\frac{k-1}{k(k+1)}z^k.
\]
It is strictly increasing there: coefficients are positive,
and termwise differentiation is valid on each compact
subinterval of \((0,1)\). Consequently the derivative of
\(G\) has the sign of
\[
       \frac{t}{1-p}-\frac{t}{p+t},
\]
which is the sign of \(t-(1-2p)\).

If \(p<1/2\), the minimum occurs at \(t=1-2p\).
Since
\[
 D(1-p\Vert p)=(1-2p)\ln((1-p)/p),
\]
the minimum is the first formula for \(g(p)\).
If \(p\geq1/2\), \(G\) increases for all \(t>0\).
The numerator and its first derivative vanish at zero,
and its second derivative at zero is \(1/[p(1-p)]\).
Taylor's formula therefore gives
\(\lim_{t\downarrow0}G(t,p)=1/[2p(1-p)]\),
the other formula. This proves the infimum statement
and the second exponential bound.

For \(p\geq1/2\), \(p(1-p)\leq1/4\) gives \(g(p)\geq2\).
For \(p<1/2\), put \(v=1-2p\in(0,1)\). Then
\[
 g(p)=\frac1v\ln\frac{1+v}{1-v}
     =\frac1v\int_0^v\frac{2}{1-s^2}\,ds\geq2.
\]
The same integral gives the continuous limit two at \(p=1/2\).
This proves the third bound.

At the endpoint \(t=1-p\), the event is that every \(X_i=1\).
Its probability is \(\prod_iP(X_i=1)\leq\prod_i\mu_i\leq p^n\).
Independent Bernoulli variables of common mean \(p\) attain this.
The divergence expression tends to \(p^n\), since
\(D(1\Vert p)=\ln(1/p)\). Larger deviations are impossible
because each \(X_i\leq1\).
If \(p=0\), every \(\mu_i=0\), forcing every \(X_i=0\) almost
surely. If \(p=1\), every \(X_i=1\) almost surely.
Applying the established result to \(1-X_i\) proves the
lower-tail assertion; an affine change proves the common
nondegenerate interval version.
\end{proof}
\end{annexitem}

\begin{annexitem}{annex:p15.conc-common-comparison}{thm:p15-conc-common-comparison}{Power-series comparisons and the reflected divergence identity}
\begin{proof}
For fixed \(p\), put \(J(t)=D(p+t\Vert p)\) on
\(-p<t<1-p\), and extend \(G(t,p)=J(t)/t^2\) continuously
at zero by \(G(0,p)=1/[2p(1-p)]\).
Twice integrating the geometric expansions of
\[
                 J''(t)=\frac1{p+t}+\frac1{1-p-t}
\]
with \(J(0)=J'(0)=0\) gives
\[
 G(t,p)=\sum_{k=0}^{\infty}
 \frac{(-1)^kp^{-(k+1)}+(1-p)^{-(k+1)}}
      {(k+1)(k+2)}\,t^k
                  \quad\bigl(|t|<\min(p,1-p)\bigr).
\]
The integrations are justified by uniform convergence on
compact subintervals. When \(0<p<1/2\), the formula
also holds at \(t=p\): the part with powers \(t/p\)
is absolutely summable there, and the other part
converges geometrically, so passage to the endpoint
is valid. This includes the source's stated series domain.

For \(p>1/2\), every positive-order coefficient is
positive, and the constant term equals \(g(p)\).
The coefficient of \(t\) is
\((2p-1)/[6p^2(1-p)^2]\).
For \(k\geq1\), the coefficient is at most
\(1/[3(1-p)^{k+1}]\); at higher orders the inequality
is strict. Consequently, on \(0<t<1-p\),
\[
 \frac{(2p-1)t}{6p^2(1-p)^2}
       <G(t,p)-g(p)
       <\frac{t}{3(1-p)(1-p-t)}.
\]
Since \(A_1/A_2=\exp\{-nt^2(G-g)\}\), reversing
the inequalities after negating proves the first pair
of ratio bounds.

At \(p=1/2\), all odd coefficients vanish, the
coefficient of \(t^2\) is \(4/3\), and the coefficient
of \(t^{2j}\), \(j\geq1\), is
\[
                       \frac{2^{2j+2}}{(2j+1)(2j+2)}.
\]
It is at most \(4^j/3\), with strict inequality for
\(j>1\). Summing the geometric upper series gives
\[
          \frac43t^2<G(t,1/2)-2<
                    \frac{4t^2}{3(1-4t^2)}.
\]
Exponentiation gives the displayed fourth-order comparisons.
The earlier proof of \(g(p)\geq2\) is strict unless
\(p=1/2\); with \(n,t>0\), this also proves the exact
equality criterion for \(A_2,A_3\).

For \(p<1/2\), let \(c=1-2p\) and \(t'=c-t\).
Then \(p+t'=1-(p+t)\in(0,1)\), even when \(t'<0\).
Direct subtraction of the two divergence expressions gives
\[
 J(t)-J(t')=(2t-c)\ln((1-p)/p)
                    =g(p)(t^2-t'^2).
\]
Thus
\[
                 J(t)-g(p)t^2=J(t')-g(p)t'^2.
\]
The elementary inequality \(\ln x\leq x-1\),
applied to both logarithms in \(D(q\Vert p)\), gives
\[
 D(q\Vert p)\leq
 q(q/p-1)+(1-q)((1-q)/(1-p)-1)
               =\frac{(q-p)^2}{p(1-p)}.
\]
Also \(\ln x>1-1/x\) for \(x>1\), by integrating
the positive derivative of \(\ln x-1+1/x\).
At \(x=(1-p)/p\), this gives \(g(p)>1/(1-p)\).
Hence
\[
 J(t)-g(p)t^2
   \leq t'^2\left(\frac1{p(1-p)}-g(p)\right)
   \leq\frac{t'^2}{p}.
\]
Together with the already proved nonnegativity of
\(J(t)-g(p)t^2\), this proves the third ratio bound.
The minimum characterization of \(G\) proves that
the ratio equals one exactly at \(t=c\).

Finally, for fixed \(p<1/2\), the same expansion at
zero gives
\[
 J(t)-g(p)t^2
    =\left(\frac1{2p(1-p)}-g(p)\right)t^2+O(t^3).
\]
If \(nt^2\to a<\infty\), then \(nt^3\to0\).
Taking the exponential of the limiting log ratio
proves the final assertion, including \(a=0\).
\end{proof}
\end{annexitem}

\begin{annexitem}{annex:p15.conc-unequal-ranges}{thm:p15-conc-unequal-ranges}{Squared ranges and both types of loss comparison}
\begin{proof}
Independence and the centered exponential lemma give
\[
 Ee^{\lambda S_n}
    =\prod_iEe^{\lambda(X_i-\mu_i)}
    \leq\exp(\lambda^2 V_n/8).
\]
For positive \(\lambda\), exponential Markov bounds the
upper tail by \(\exp(-\lambda t+\lambda^2V_n/8)\).
When \(V_n>0\), its minimum is at \(\lambda=4t/V_n\)
and equals \(\exp(-2t^2/V_n)\). Replacing all centered
variables by their negatives gives the lower tail.
The two events at a positive threshold are disjoint, and
addition gives the two-sided bound.
When \(V_n=0\), every centered variable is zero almost surely,
so positive deviations are impossible.

For a real weight \(w_i\), the interval of \(w_iX_i\)
has length \(|w_i|(b_i-a_i)\), including negative or zero
weights. Applying the preceding result to those independent
variables proves the weighted formula.
For the difference of two independent sample means, use
weights \(1/m\) on the \(Y\)'s and \(-1/n\) on the \(Z\)'s.
The sum of squared ranges is
\((b-a)^2(m/m^2+n/n^2)\), as stated.

For paired losses \(A_i,B_i\in[0,1]\) that are functions of
the same independent observation \(Z_i\), define \(W_i=A_i-B_i\).
The \(W_i\)'s are independent across records even when
\(A_i,B_i\) are dependent within a record.
Each lies in \([-1,1]\). The weighted mean uses squared
range sum \(n(2/n)^2=4/n\), giving
\(\exp(-nt^2/2)\).
The two-sample independence hypothesis was not used to
justify this separate paired calculation.
\end{proof}
\end{annexitem}

\begin{annexitem}{annex:p15.conc-martingale}{thm:p15-conc-martingale}{Conditional exponential bounds and maximal deviations}
\begin{proof}
Apply the deterministic chord inequality to \(D_i\)
and take conditional expectation given \(\mathcal F_{i-1}\).
Its conditional mean is zero, so the same two-point
calculation as in the exponential lemma yields
\[
 E(e^{\lambda D_i}\mid\mathcal F_{i-1})
                  \leq e^{\lambda^2(b_i-a_i)^2/8}.
\]
The zero-width case gives \(D_i=0\) almost surely.
All conditional quantities are integrable because their
variables have deterministic finite bounds.
Multiplying by \(e^{\lambda M_{i-1}}\), taking expectation
and iterating prove
\(Ee^{\lambda M_n}\leq e^{\lambda^2V_n/8}\).

For fixed real \(\lambda\), conditional Jensen also gives
\(E(e^{\lambda D_i}\mid\mathcal F_{i-1})\geq1\).
Thus \(e^{\lambda M_k}\) is a nonnegative submartingale.
For \(\lambda>0\), Doob's finite-horizon inequality
\cref{thm:p15-doob-weak-maximal} gives
\[
 P\{\max_{k\leq n}M_k\geq t\}
   =P\{\max_{k\leq n}e^{\lambda M_k}\geq e^{\lambda t}\}
   \leq e^{-\lambda t}Ee^{\lambda M_n}.
\]
Optimizing as before proves the upper maximum bound.
Apply the same argument to \(-M\), then use the union
bound for the absolute maximum. The events in this
last union need not be disjoint. If \(V_n=0\), all
increments are zero and all maxima of positive size
are impossible.

For the source's common-range claim, the conditional chord
bound is
\[
 E(e^{\lambda X_i}\mid\mathcal F_{i-1})
                 \leq1-\mu_i+\mu_i e^\lambda.
\]
The \(\mu_i\)'s are deterministic, so iteration yields
the same product as in the independent common-range proof.
The centered partial sums are martingales; their
exponentials are nonnegative submartingales by conditional
Jensen. Doob's inequality and the same arithmetic-geometric
mean step give the optimized divergence bound for the
maximum. The two algebraic refinements depend only on
\(p,t\) and therefore follow unchanged.
\end{proof}
\end{annexitem}

\begin{annexitem}{annex:p15.conc-benchmarks}{thm:p15-conc-benchmarks}{Moment bounds, Gaussian tails and the nonnegative refinement}
\begin{proof}
Markov's inequality for \(W^2\) gives the two-sided
variance bound. For the one-sided bound, apply
\cref{thm:p15-cantelli-maximal} to the one-step martingale
with initial value zero and terminal value \(W\);
its filtration is trivial initially and contains \(W\)
at the terminal step. This gives \(v/(v+u^2)\).
If \(v=0\), nonnegativity and zero expectation of
\(W^2\) imply \(W=0\) almost surely.

Completing the square in the normalized Gaussian density
gives \(Ee^{\lambda Z}=e^{\lambda^2/2}\).
Exponential Markov optimized at \(\lambda=y>0\)
gives \(P(Z\geq y)\leq e^{-y^2/2}\).
Let \(I(y)=\int_y^\infty\phi(t)\,dt\).
Using \(\phi'(t)=-t\phi(t)\), integration by parts gives
\[
 I(y)=\frac{\phi(y)}y-\int_y^\infty\frac{\phi(t)}{t^2}\,dt.
\]
The boundary term at infinity vanishes because
\(\phi(t)/t\to0\). The remaining integral is nonnegative
and at most \(I(y)/y^2\). Rearranging supplies both
displayed tail bounds. Symmetry gives \(I(y)=\Phi(-y)\),
and multiplying by \(y/\phi(y)\) squeezes the ratio
between \(y^2/(1+y^2)\) and one, proving its limit.

The row assumptions are the Lyapunov condition of
\cref{thm:p15-clt-lyapunov} with exponent parameter one,
so the normalized centered sum converges weakly to a
standard normal variable. That distribution has a
continuous CDF at every real \(y\). Squeezing an
inclusive upper-tail event between the corresponding
strict events at \(y-\varepsilon\) and \(y+\varepsilon\),
then sending \(\varepsilon\downarrow0\), proves the
inclusive tail limit. The extra variance-divergence
condition is retained from the original displayed
benchmark; Lyapunov's normalized theorem does not
require inferring it from an unrelated source.

For the unbounded-range example, take \(X_1\) to equal
\(n(\mu+t)\) with probability \(\mu/(\mu+t)\) and zero
otherwise; let \(X_2,\ldots,X_n\) be identically zero.
These variables are independent and their average
mean is \(\mu\). The event \(\bar X-\mu\geq t\)
occurs exactly at the positive value of \(X_1\), so
its probability attains the nonnegative Markov bound.

Finally, put \(U=\sum_iX_i\geq0\) and \(a=n(p+t)>0\).
For \(\lambda>0\),
\[
       \mathbf1_{\{U\geq a\}}
            \leq\frac{e^{\lambda U}-1}{e^{\lambda a}-1}.
\]
The right side is nonnegative below the threshold
and at least one above it. Taking expectation and
using the earlier chord/product/arithmetic-mean bound
gives the displayed refinement.
At \(\lambda_*\), let \(M=(1-p+pe^{\lambda_*})^n\)
and \(C=e^{\lambda_*n(p+t)}\). Then
\(1<M<C\), since \(M/C=e^{-nD(p+t\Vert p)}<1\).
The strict positivity of the divergence follows
from its already proved lower bound \(2t^2\).
Direct cross multiplication shows
\((M-1)/(C-1)<M/C\), completing the comparison.
\end{proof}
\end{annexitem}

\begin{annexitem}{annex:p15.conc-finite-learning}{thm:p15-conc-finite-learning}{Uniform control before selection}
\begin{proof}
For fixed \(h\), the \(n\) losses are independent and
bounded in \([0,1]\). The sum inequality at threshold
\(nb\) gives
\[
 P\{|\widehat R_n(h)-R(h)|>b\}\leq2e^{-2nb^2}.
\]
Taking a union over the \(r\) fixed hypotheses gives
probability at most \(2r e^{-2nb^2}\).
At the stated \(b_n\) this is \(\delta\).
The event thus controls every hypothesis simultaneously.
Independence across the different hypotheses is unnecessary.
Measurability follows from finite sums and finite unions.

On that event, for every comparison hypothesis \(h\)
and every \(\eta\)-approximate minimizer,
\[
 \begin{split}
 R(\widehat h)
 &\leq\widehat R_n(\widehat h)+b_n\\
 &\leq\widehat R_n(h)+\eta+b_n\\
 &\leq R(h)+2b_n+\eta.
 \end{split}
\]
Minimizing the right side over the finite class gives
the risk guarantee. A tie-resolved exact minimizer
exists and is measurable: its selection events are
finite comparisons of measurable empirical risks.
The condition \(2b_n\leq\varepsilon\) is exactly
\(n\geq2\ln(2r/\delta)/\varepsilon^2\).
Finally, on the same event, every true risk is in
its untruncated interval and also in \([0,1]\).
It is therefore in the displayed clipped interval.
This applies to any selected member because the
coverage event already quantifies over all of them.
\end{proof}
\end{annexitem}

\begin{annexitem}{annex:p15.conc-weighted-learning}{thm:p15-conc-weighted-learning}{Countable allocation and penalized selection}
\begin{proof}
For each fixed \(h\), the previous bound at radius \(b_h\)
is \(\delta w_h\). Countable subadditivity bounds the
probability of any violation by
\(\delta\sum_hw_h\leq\delta\).
On its complement, for every \(h\),
\[
 \begin{split}
 R(\widehat h)
 &\leq\widehat R_n(\widehat h)+b_{\widehat h}\\
 &\leq\widehat R_n(h)+b_h+\eta\\
 &\leq R(h)+2b_h+\eta.
 \end{split}
\]
Taking the infimum proves the conclusion, without
requiring that the comparison infimum be attained.

For existence of an approximate selection, enumerate
the countable class and let \(J\) be the infimum of
its empirical penalized scores. These scores are
nonnegative and finite, so \(0\leq J<\infty\).
Their countable infimum is measurable.
For \(\eta>0\), some score is less than \(J+\eta\).
Choose the first such index. Its event is the
intersection of the corresponding strict comparison
with finitely many failed earlier comparisons,
and is measurable. This proves existence.
An infimum at \(\eta=0\) need not be a minimum.
None of these measurability statements asserts an
effective algorithm for computing an arbitrary
countable infimum.
\end{proof}
\end{annexitem}

\begin{annexitem}{annex:p15.conc-validation}{thm:p15-conc-validation}{Conditioning on the entire training object}
\begin{proof}
Fix a training value \(t\). The \(K\) candidate loss
functions \(\ell_j(t,\cdot)\) are now fixed bounded
measurable functions. The validation data have the
product law \(P^{\otimes m}\), independent of this value.
Applying the finite-class bound gives conditional
failure at most \(\delta\), uniformly in \(t\).
The risks \(R_j(t)\) are measurable by nonnegative
integration, and the finite union of failure events
is product measurable by the assumed joint loss
measurability. Integrating its probability against
the training law preserves the bound \(\delta\).
This is a product-measure argument and does not
require a pointwise conditional probability on a
null training event.

On the resulting uniform event, comparison of the
validation minimizer with any of the \(K\) candidates
uses exactly the three inequalities in the empirical
risk proof. The excess risk is at most twice the
validation radius. If only one trained candidate
is evaluated, the same proof has \(K=1\).
If the candidates depend on the validation values,
fixing \(t\) no longer fixes their loss functions.
This proof then supplies no such \(K\)-candidate bound.
\end{proof}
\end{annexitem}

\begin{annexitem}{annex:p15.conc-noisy-example}{thm:p15-conc-noisy-example}{All four risks and the specified sample}
\begin{proof}
The joint probabilities of \((X,Y)=(0,0),(0,1),(1,0),(1,1)\)
are \(3/8,1/8,1/8,3/8\). Thus the constant-zero and
constant-one errors are each \(1/2\).
The identity predictor errs only at \((0,1),(1,0)\),
with total mass \(1/4\); its complement errs on the
other two pairs, with mass \(3/4\).
In particular none of these four has zero risk.

The ordered sample's probability is
\((3/8)(1/8)(3/8)(3/8)=27/4096\).
The constant-zero predictor makes three errors,
the constant-one predictor one, the identity one,
and the complement three. The first empirical
minimizer in the displayed order is therefore
constant one. Its risk exceeds the best risk
by \(1/2-1/4=1/4\).

For \(r=4,\delta=1/20,\varepsilon=1/10\), the
sufficient bound is \(2\ln160/(1/10)^2=200\ln160\).
For example, the convergent identity
\[
 \ln x=2\sum_{j=0}^{\infty}
       \frac{z^{2j+1}}{2j+1},\qquad z=\frac{x-1}{x+1},
\]
with a geometric bound on its positive remainder,
applied to \(\ln160=7\ln2+\ln(5/4)\), gives
\(5.07517<\ln160<5.07518\).
Hence \(1015<200\ln160<1016\), proving the stated ceiling.
It is a sufficient guarantee from the uniform bound;
the example does not assert a matching lower bound.
\end{proof}
\end{annexitem}

\begin{annexitem}{annex:p15.conc-selection-example}{thm:p15-conc-selection-example}{The exact multiple-search probability}
\begin{proof}
Write the observed feature matrix as \(B_{ij}\),
with \(1\leq i\leq n,\ 1\leq j\leq K\).
Its \(nK\) entries are independent fair bits.
For each fixed \(j\), true classification risk is
\(P(B_{1j}=1)=1/2\). Zero empirical error means
that every bit in column \(j\) is zero, an event
of probability \(2^{-n}\).
Different columns involve disjoint sets of independent
bits, so their all-zero events are mutually independent.
The probability that no column is all zero is
\((1-2^{-n})^K\). Taking its complement proves the formula.
Every selected column still represents a predictor
of true risk \(1/2\); its sample fit does not alter
that fixed risk. The union over candidates is
therefore a necessary part of the sampling question.
\end{proof}
\end{annexitem}

\begin{annexitem}{annex:p15.conc-dependence-boundaries}{thm:p15-conc-dependence-boundaries}{Fifteen pairwise-independent bits and a growing range}
\begin{proof}
Use four independent fair bits \(Z=(Z_1,\ldots,Z_4)\),
viewed in the vector space over \(\mathbb F_2=\{0,1\}\)
with arithmetic modulo two. For each nonzero vector
\(a\in\mathbb F_2^4\), set \(X_a=\mathbf1_{\{a\cdot Z=0\}}\).
There are fifteen such variables.
Each nonzero linear parity is a fair bit, since
flipping a coordinate with coefficient one pairs its
zero and one inputs.
Distinct nonzero vectors \(a,b\) are linearly independent
over this two-element field. The map
\(Z\mapsto(a\cdot Z,b\cdot Z)\) has four equally sized
fibers: two independent coefficient columns determine
two coordinates uniquely for any requested output,
while the other two coordinates are free.
Thus \(X_a,X_b\) are independent fair binary variables.

At \(Z=0\), all fifteen \(X_a\)'s are one.
At any nonzero \(Z\), exactly eight of the sixteen
vectors \(a\), including the zero vector, satisfy
\(a\cdot Z=0\), so exactly seven nonzero vectors do.
The centered sum is consequently \(15/2\) at \(Z=0\)
and \(-1/2\) otherwise. Its upper event at \(15/2\)
has probability \(1/16\).
The independent range bound would be \(e^{-15/2}\).
Since \(e^{15/2}>1+15/2+(15/2)^2/2>16\), that
purported bound is strictly smaller than the actual
probability. Pairwise independence is insufficient.

For the range example, at each fixed \(n\) let
\(Y=2n\) with probability \(1/(2n)\), and zero otherwise.
Its expectation is one. In an independent sample
of size \(n\), every observation is zero with probability
\((1-1/(2n))^n\geq1/2\). The last inequality follows
from \((1-u)^n\geq1-nu\) for \(0\leq u\leq1\),
proved by induction after multiplying by \(1-u\).
Every member of this family has a finite mean, but
its range length is \(2n\). No common fixed range
was supplied, so the range-free exponential conclusion
does not follow.
\end{proof}
\end{annexitem}

\begin{annexitem}{annex:p15.conc-time-uniform}{thm:p15-conc-time-uniform}{A countable probability allocation over observation times}
\begin{proof}
For \(V_n>0\), substitute
\[
          t_n=\sqrt{(V_n/2)\ln(2n(n+1)/\delta)}
\]
in the two-sided martingale bound for the fixed time \(n\).
The probability of \(|M_n|>t_n\) is at most
\(\delta/[n(n+1)]\). If \(V_n=0\), all increments
through \(n\) are zero and this strict-violation event
is empty for radius zero.
Countable subadditivity and the telescoping sum
\[
       \sum_{n=1}^{\infty}\frac{1}{n(n+1)}
          =\sum_{n=1}^{\infty}\left(\frac1n-\frac1{n+1}\right)=1
\]
give probability at most \(\delta\) of any violation.
On its complement every displayed finite-time bound
holds at once.

For independent \(X_i\in[0,1]\) with common mean \(\mu\),
the centered increments have bounds \(-\mu,1-\mu\)
and width one. Thus \(V_n=n\) and
\(M_n=n(\bar X_n-\mu)\). Division gives the mean band.
For an almost surely finite random time \(\tau\),
intersect the simultaneous event with \(\{\tau<\infty\}\),
an event of probability one, and substitute \(n=\tau\).
No assumption that the stopping time is independent
of the observations is needed for this substitution.
\end{proof}
\end{annexitem}
